\chapter{Mehrfaches Bilden von Paarwerten}
\label{b14guide:ch:iteration}

\section{Lokale Assoziativität}

\begin{theorem}[Umklammern und Existenz der äußeren Paarwerte]
\label{b14guide:thm:associative}
Seien $x\wedge y$ und $y\wedge z$ vorhanden. Dann gilt
\[
 \InfEx(x\wedge y,z)\Longleftrightarrow\InfEx(x,y\wedge z).
\]
Wenn eine dieser Bedingungen erfüllt ist, sind beide Seiten
der folgenden Gleichung definiert und gleich:
\[
 (x\wedge y)\wedge z=x\wedge(y\wedge z).
\]
Seien stattdessen $x\vee y$ und $y\vee z$ vorhanden. Dann gilt
\[
 \SupEx(x\vee y,z)\Longleftrightarrow\SupEx(x,y\vee z),
\]
und bei Erfüllung einer dieser Bedingungen
\[
 (x\vee y)\vee z=x\vee(y\vee z).
\]
Die Existenz beider inneren Paarwerte bleibt jeweils Voraussetzung.
\end{theorem}

\begin{theorem}[Auswerten eines dreifachen Ausdrucks]
\label{b14guide:thm:triple}
Sind $x\wedge y$ und $(x\wedge y)\wedge z$ vorhanden, so gilt
für jedes $u\in P$
\[
 u\leq(x\wedge y)\wedge z
 \Longleftrightarrow (u\leq x\land u\leq y\land u\leq z).
\]
Sind $x\vee y$ und $(x\vee y)\vee z$ vorhanden, so gilt
\[
 (x\vee y)\vee z\leq u
 \Longleftrightarrow (x\leq u\land y\leq u\land z\leq u).
\]
Die Existenz von $y\wedge z$ beziehungsweise $y\vee z$
wird für diese Auswertung nicht benötigt.
\end{theorem}

\section{Endliche Ausdrücke aus einer Paaroperation}

\begin{definition}[Zulässiger Paarterm]
\label{b14guide:def:term}
Ein \emph{endlicher Infimumspaarterm} entsteht aus endlich vielen,
mindestens einem, Elementen von $P$ durch vollständiges
Klammern mit $\wedge$. Er heißt \emph{zulässig}, wenn alle
inneren und äußeren Paarwerte bei der Auswertung existieren.
Ein einzelnes Element ist bereits ein zulässiger Term.
Für \emph{Supremumspaarterme} gilt dieselbe Definition mit $\vee$.
Ein Term verwendet hier stets nur eine der beiden Operationen.
\end{definition}

\begin{theorem}[Universelle Eigenschaft eines zulässigen Paarterms]
\label{b14guide:thm:term}
Sei $t$ der Wert eines zulässigen Infimumspaarterms mit
Operanden $x_1,\ldots,x_n$, wobei $n\geq1$. Dann gilt
\[
 \forall u\in P\quad
 \bigl(u\leq t\Longleftrightarrow
       (u\leq x_1\land\cdots\land u\leq x_n)\bigr).
\]
Für den Wert $s$ eines zulässigen Supremumspaarterms gilt
\[
 \forall u\in P\quad
 \bigl(s\leq u\Longleftrightarrow
       (x_1\leq u\land\cdots\land x_n\leq u)\bigr).
\]
\end{theorem}

\begin{theorem}[Unabhängigkeit zulässiger Auswertungen]
\label{b14guide:thm:term-equality}
Zwei zulässige Infimumspaarterme haben denselben Wert,
wenn in ihnen dieselben Elemente vorkommen. Reihenfolge,
Klammerung und Wiederholungen ändern den Wert nicht.
Dasselbe gilt für zwei zulässige Supremumspaarterme.

Die Aussage setzt die Zulässigkeit \emph{beider} Terme voraus.
Aus der Zulässigkeit einer Klammerung folgt im allgemeinen
nicht die Zulässigkeit jeder anderen Klammerung.
\end{theorem}

\begin{theorem}[Vertauschungsregel für vier Operanden]
\label{b14guide:thm:four}
Sind sämtliche Teilwerte auf beiden Seiten vorhanden, so gilt
\[
\begin{aligned}
 (x\wedge y)\wedge(u\wedge v)
   &=(x\wedge u)\wedge(y\wedge v),\\
 (x\vee y)\vee(u\vee v)
   &=(x\vee u)\vee(y\vee v).
\end{aligned}
\]
Für die erste Gleichung werden nur Paarinfima, für die zweite
nur Paarsuprema vorausgesetzt.
\end{theorem}